\section{Elementos flutuantes e conteúdo técnico}

\subsection{Figura}

Use \texttt{figure} para que o \LaTeX\ posicione a imagem, numere a legenda e
inclua o item na lista de figuras. O marcador de \verb|\label| deve vir depois de
\verb|\caption|.

\begin{lstlisting}[language={[LaTeX]TeX},caption={Inserção de uma figura}]
\begin{figure}[H]
  \centering
  \includegraphics[width=.35\linewidth]{images/unesp.png}
  \caption{Logotipo utilizado no modelo}
  \label{fig:unesp}
  \fonteimagem{Fonte: arquivo do modelo.}
\end{figure}
\end{lstlisting}

\begin{figure}[H]
  \centering
  \includegraphics[width=.35\linewidth]{images/unesp.png}
  \caption{Logotipo utilizado no modelo}
  \label{fig:unesp}
  \fonteimagem{Fonte: arquivo do modelo.}
\end{figure}

\subsection{Subfiguras}

\begin{lstlisting}[language={[LaTeX]TeX},caption={Duas imagens lado a lado}]
\begin{figure}[H]
  \centering
  \begin{subfigure}{.3\linewidth}
    \includegraphics[width=\linewidth]{images/unesp.png}
    \caption{Primeira imagem}
  \end{subfigure}\hfill
  \begin{subfigure}{.3\linewidth}
    \includegraphics[width=\linewidth]{images/unesp.png}
    \caption{Segunda imagem}
  \end{subfigure}
  \caption{Exemplo com subfiguras}
\end{figure}
\end{lstlisting}

\begin{figure}[H]
  \centering
  \begin{subfigure}{.3\linewidth}
    \includegraphics[width=\linewidth]{images/unesp.png}
    \caption{Primeira imagem}
  \end{subfigure}\hfill
  \begin{subfigure}{.3\linewidth}
    \includegraphics[width=\linewidth]{images/unesp.png}
    \caption{Segunda imagem}
  \end{subfigure}
  \caption{Exemplo com subfiguras}
  \label{fig:subfiguras}
\end{figure}

\subsection{Tabela}

A coluna \texttt{X} do ambiente \texttt{tabularx} distribui o espaço disponível.
\verb|\multirow| permite que uma célula ocupe mais de uma linha.

\begin{lstlisting}[language={[LaTeX]TeX},caption={Tabela com largura ajustável}]
\begin{table}[H]
  \centering
  \caption{Exemplo de tabela}
  \label{tab:exemplo}
  \begin{tabularx}{\linewidth}{|l|X|c|}
    \hline
    Grupo & Descricao & Valor \\ \hline
    \multirow{2}{*}{A} & Primeiro resultado & 10 \\ \cline{2-3}
                       & Segundo resultado  & 20 \\ \hline
  \end{tabularx}
\end{table}
\end{lstlisting}

\begin{table}[H]
  \centering
  \caption{Exemplo de tabela}
  \label{tab:exemplo}
  \begin{tabularx}{\linewidth}{|l|X|c|}
    \hline
    Grupo & Descrição & Valor \\ \hline
    \multirow{2}{*}{A} & Primeiro resultado & 10 \\ \cline{2-3}
                       & Segundo resultado  & 20 \\ \hline
  \end{tabularx}
\end{table}

\subsection{Equações}

Equações entre \verb|\[| e \verb|\]| não são numeradas. Para criar número e
referência, use o ambiente \texttt{equation}.

\begin{lstlisting}[language={[LaTeX]TeX},caption={Equação numerada}]
\begin{equation}
  E = mc^2
  \label{eq:energia}
\end{equation}
\end{lstlisting}

\begin{equation}
  E = mc^2
  \label{eq:energia}
\end{equation}
